\documentclass[12pt,a4paper]{report}
\usepackage[utf8]{inputenc}
\usepackage[T1]{fontenc}
\usepackage[spanish,es-nodecimaldot,es-tabla,es-noquoting]{babel}

% Paquetes de graficos y diagramas
\usepackage{graphicx}
\graphicspath{{../img/}{img/}}
\usepackage{tikz}
\usetikzlibrary{shapes.geometric,shapes.misc,arrows.meta,positioning,calc,fit,backgrounds,shadows.blur}

% Tipografia Times New Roman estandar en 12pt
\usepackage{newtxtext,newtxmath}

% Margenes estandar normalizados (3cm izq, 2.5cm der/sup/inf)
\usepackage{geometry}
\geometry{
  top=2.5cm,
  bottom=2.5cm,
  left=3.0cm,
  right=2.5cm,
  headheight=15pt,
  headsep=20pt,
  footskip=25pt
}

% Espaciado e interlineado
\usepackage{setspace}
\setstretch{1.15}
\setlength{\parskip}{6pt}
\setlength{\parindent}{0pt}

% Vinetas redondas formales
\usepackage{enumitem}
\setlist[itemize,1]{label=\textbullet, leftmargin=1.5em, itemsep=2pt, parsep=0pt}
\setlist[itemize,2]{label=\normalfont\bfseries\textperiodcentered, leftmargin=1.5em, itemsep=2pt, parsep=0pt}
\setlist[enumerate]{leftmargin=1.5em, itemsep=2pt, parsep=0pt}

% Color del texto y paleta institucional
\usepackage{xcolor}
\definecolor{espochBlue}{RGB}{0,51,102}
\definecolor{espochRed}{RGB}{192,0,0}
\definecolor{techGray}{RGB}{245,247,250}
\definecolor{borderGray}{RGB}{209,213,219}

% Tablas profesionales segun estandares ISO/APA
\usepackage{booktabs}
\usepackage{tabularx}
\usepackage{longtable}
\usepackage{array}
\usepackage{caption}
\captionsetup[table]{position=top,skip=6pt,font={small,bf},labelfont=bf}
\captionsetup[figure]{position=bottom,skip=6pt,font={small,bf},labelfont=bf}

% Listados de codigo preformateado
\usepackage{listings}
\lstset{
  basicstyle=\small\ttfamily,
  frame=single,
  rulecolor=\color{black},
  framerule=0.5pt,
  breaklines=true,
  numbers=none,
  backgroundcolor=\color{white},
  tabsize=2,
  inputencoding=utf8,
  extendedchars=true,
  literate=
    {á}{{\'a}}1 {é}{{\'e}}1 {í}{{\'{\i}}}1 {ó}{{\'o}}1 {ú}{{\'u}}1
    {Á}{{\'A}}1 {É}{{\'E}}1 {Í}{{\'I}}1 {Ó}{{\'O}}1 {Ú}{{\'U}}1
    {ñ}{{\~n}}1 {Ñ}{{\~N}}1 {ü}{{\"u}}1 {Ü}{{\"U}}1
}

% Encabezados y pies de pagina formales
\usepackage{fancyhdr}
\pagestyle{fancy}
\fancyhf{}
\fancyhead[L]{\small\nouppercase{\leftmark}}
\fancyhead[R]{\small ESPOCH - Metáfora y Stack Tecnológico}
\fancyfoot[C]{\small\thepage}
\renewcommand{\headrulewidth}{0.4pt}
\renewcommand{\footrulewidth}{0.4pt}

% Formato de titulos
\usepackage{titlesec}
\titleformat{\chapter}[display]
  {\normalfont\Large\bfseries}{\chaptertitlename\ \thechapter}{10pt}{\huge\bfseries}
\titlespacing*{\chapter}{0pt}{-10pt}{15pt}

\titleformat{\section}
  {\normalfont\large\bfseries}{\thesection}{1em}{}
\titlespacing*{\section}{0pt}{12pt}{4pt}

\titleformat{\subsection}
  {\normalfont\normalsize\bfseries}{\thesubsection}{1em}{}
\titlespacing*{\subsection}{0pt}{10pt}{3pt}

\titleformat{\subsubsection}
  {\normalfont\normalsize\bfseries}{\thesubsubsection}{1em}{}
\titlespacing*{\subsubsection}{0pt}{8pt}{2pt}

% Enlaces sin color pero navegables
\usepackage{hyperref}
\hypersetup{
  hidelinks=true,
  pdfauthor={Equipo de Ingenieria de Desarrollo y Aseguramiento Tecnico - ESPOCH},
  pdftitle={Documento de Metáfora y Stack Tecnológico - ESPOCH},
  pdfsubject={Proyecto Integrador de Saberes - Quinto Semestre}
}

\begin{document}

% PORTADA ACADEMICA INSTITUCIONAL
\begin{titlepage}
  \centering
  \thispagestyle{empty}
  \vspace*{0.3cm}

  \includegraphics[width=3.2cm]{LogoESPOCH.png}\\[0.4cm]

  {\normalsize\bfseries ESCUELA SUPERIOR POLITÉCNICA DE CHIMBORAZO}\\[0.2cm]
  {\small\bfseries FACULTAD DE INFORMÁTICA Y ELECTRÓNICA}\\[0.1cm]
  {\small\bfseries CARRERA DE TECNOLOGÍAS DE LA INFORMACIÓN}\\[1.0cm]

  {\large\bfseries SISTEMA WEB DE MONITOREO Y GESTIÓN\\[0.15cm]
  DE LA INFRAESTRUCTURA DE REDES - ESPOCH}\\[0.7cm]

  {\normalsize\bfseries DOCUMENTO DE METÁFORA, ARQUITECTURA DIRIGIDA POR EVENTOS\\[0.15cm]
  Y SELECCIÓN DEL STACK TECNOLÓGICO}\\[0.15cm]
  {\small Arquitectura Reactiva en Tiempo Real, Ecosistema de Herramientas y Justificación Técnica}\\[1.0cm]

  \begin{minipage}[t]{0.45\textwidth}
    {\small\textbf{Roles de Ingeniería:}}\\[0.25cm]
    \begin{itemize}[leftmargin=1.2em,itemsep=2pt,parsep=0pt]
      \small
      \item Líder Técnico y Arquitecto Backend
      \item Desarrollador Frontend y Diseñador UI/UX
      \item Especialista en Bases de Datos y QA Lead
      \item Ingeniero DevOps y Redes
    \end{itemize}
  \end{minipage}
  \hfill
  \begin{minipage}[t]{0.50\textwidth}
    {\small\textbf{Asignaturas y Docentes Responsables:}}\\[0.25cm]
    \begin{itemize}[leftmargin=1.2em,itemsep=2pt,parsep=0pt]
      \small
      \item Infraestructura TI: Ing. Tatiana Landivar
      \item Tecnología Web: Ing. Pedro Aguilar
      \item Conmutación y Enrutamiento: Ing. Juan Pazmiño
      \item Big Data: Ing. Geovanny Silva
      \item Diseño y Multimedia: Ing. Ramiro Bonilla
    \end{itemize}
  \end{minipage}

  \vfill

  {\small\textbf{Período Académico:} Quinto Semestre (PAO 5)}\\[0.05cm]
  {\small\textbf{Documento Oficial de Ingeniería de Software (Versión 2.0 - Arquitectura EDA)}}\\[0.6cm]

  {\small Orellana - Ecuador}\\[0.15cm]
  {\small Noviembre de 2026}
\end{titlepage}

\tableofcontents
\newpage

\listoffigures
\newpage

\listoftables
\newpage

\chapter{Fundamentos y Metáfora del Sistema}

\section{Metáfora Conceptual: Centro de Operaciones de Red (NOC)}
La metáfora del sistema proporciona un modelo mental claro y unificado del funcionamiento de la plataforma:
\begin{itemize}
  \item \textbf{Consola de Visualización (Frontend):} Proyecta en pantallas interactivas el mapa de topología, indicadores clave de desempeño (KPIs) y el panel de alertas acústico-visuales.
  \item \textbf{Canal de Telemetría (WebSockets):} Enlace bidireccional continuo que transmite métricas y notificaciones en tiempo real sin requerir peticiones periódicas.
  \item \textbf{Centro de Control y Despacho (Backend API y Bus de Eventos):} Procesa peticiones administrativas seguras (autenticación JWT, gestión de usuarios, configuración) y gestiona el flujo asíncrono de eventos internos.
  \item \textbf{Motor de Sondeo (Servicio de Monitoreo):} Servicio autónomo en segundo plano que ejecuta comprobaciones periódicas de Capa 2 (ARP) y Capa 3 (ICMP ping) sobre los equipos registrados.
  \item \textbf{Almacén Central (MongoDB):} Repositorio documental que guarda inventario, métricas agregadas y bitácora histórica de eventos con políticas de retención automática (TTL).
\end{itemize}

\begin{figure}[htbp]
  \centering
  \begin{tikzpicture}[
    scale=0.9, every node/.style={transform shape},
    box/.style={rectangle, draw=black, thick, fill=white, rounded corners=3pt, align=center, minimum height=1.2cm, minimum width=3.4cm, font=\small},
    daemon/.style={rectangle, draw=espochRed, thick, fill=techGray, rounded corners=3pt, align=center, minimum height=1.2cm, minimum width=3.4cm, font=\small},
    bus/.style={rectangle, draw=espochBlue, very thick, fill=techGray, rounded corners=4pt, align=center, minimum height=1.0cm, minimum width=4.2cm, font=\small\bfseries},
    db/.style={cylinder, draw=espochBlue, thick, fill=techGray, shape border rotate=90, aspect=0.25, minimum height=1.3cm, minimum width=3.0cm, align=center, font=\small},
    arrow/.style={-{Stealth[scale=1.0]}, thick},
    biarrow/.style={{Stealth[scale=1.0]}-{Stealth[scale=1.0]}, thick}
  ]
    % Nodos
    \node[box] (ui) at (0,4.2) {\textbf{Consola NOC}\\[2pt]Frontend React 18};
    \node[box] (api) at (0,2.1) {\textbf{Servidor de Aplicación}\\[2pt]API REST Express};
    \node[bus] (bus) at (0,0.5) {\textbf{Bus de Eventos Interno}\\[2pt]\footnotesize \texttt{EventEmitter} en Node.js};
    \node[daemon] (daemon) at (-4.5,0.5) {\textbf{Motor de Sondeo}\\[2pt]ICMP Ping / ARP-Scan};
    \node[db] (db) at (4.5,0.5) {\textbf{Base de Datos}\\[2pt]MongoDB 7.0 (TTL)};
    \node[box, fill=techGray] (lan) at (-4.5,-1.8) {\textbf{Red Local (LAN)}\\[2pt]Switches, Routers, PCs};

    % Conexiones
    \draw[arrow] (ui.south west) -- node[left, font=\scriptsize] {HTTP REST} (api.north west);
    \draw[biarrow] (ui.south east) -- node[right, font=\scriptsize] {WebSockets (Tiempo Real)} (api.north east);
    \draw[arrow] (api) -- (bus);
    \draw[arrow] (daemon) -- node[above, font=\scriptsize] {Emite Eventos} (bus);
    \draw[arrow] (bus) -- node[above, font=\scriptsize] {Persistencia} (db);
    \draw[arrow] (daemon) -- node[left, font=\scriptsize] {Sondeos ICMP/ARP} (lan);
    \draw[arrow] (lan) -- node[right, font=\scriptsize] {Respuestas} (daemon);

  \end{tikzpicture}
  \caption{Diagrama de interacción general de los componentes del sistema}\label{fig:arquitectura_general}
\end{figure}

\newpage
\chapter{Arquitectura Dirigida por Eventos (EDA)}

\section{Estructura del Flujo Reactivo}
El monitoreo de red opera sobre condiciones asíncronas: caídas de enlace, fluctuaciones de latencia e incorporación de nuevos hosts. La Arquitectura Dirigida por Eventos desacopla la detección de anomalías de su procesamiento y visualización, garantizando que el sistema reaccione en milisegundos sin bloquear el servidor web.

\begin{figure}[htbp]
\centering
\begin{tikzpicture}[
  scale=0.88, every node/.style={transform shape},
  arrow/.style={-{Stealth[scale=0.85]}, thick},
  actor/.style={circle, draw=black, thick, fill=white, minimum size=0.42cm},
  comp/.style={rectangle, draw=black, thick, fill=techGray, rounded corners=2pt, minimum height=1.05cm, minimum width=2.45cm, align=center, font=\scriptsize},
  broker/.style={cylinder, draw=espochBlue, very thick, fill=white, shape border rotate=0, aspect=0.12, minimum height=1.0cm, minimum width=6.2cm, align=center, font=\small\bfseries}
]

  % Actor
  \node[actor] (head) at (0,7.1) {};
  \draw[thick] (0,6.89) -- (0,6.25); % cuerpo
  \draw[thick] (-0.3,6.6) -- (0.3,6.6); % brazos
  \draw[thick] (0,6.25) -- (-0.25,5.65); % pierna izq
  \draw[thick] (0,6.25) -- (0.25,5.65); % pierna der
  \node[font=\footnotesize\bfseries] at (0,5.35) {Operador / Administrador NOC};

  % Frontend Web
  \node[comp, minimum width=3.4cm] (front) at (0,4.1) {\textbf{Frontend Web}\\[1pt]React 18 + Socket.io Client};

  % Modulo Central / Controlador API
  \node[comp, minimum width=3.4cm] (api) at (0,2.4) {\textbf{ModuloControlador}\\[1pt]API REST Express / Auth};

  % Broker / Bus Central
  \node[broker] (bus) at (0,0.8) {Bus de Eventos Central\\(\texttt{EventEmitter} en Node.js)};

  % Modulos Inferiores
  \node[comp] (m1) at (-5.6,-1.7) {\textbf{ModuloMonitoreo}\\[1pt]Ping / ARP-Scan};
  \node[comp] (m2) at (-2.8,-1.7) {\textbf{ModuloAlertas}\\[1pt]Reglas de Umbrales};
  \node[comp] (m3) at (0.0,-1.7) {\textbf{ModuloPersistencia}\\[1pt]MongoDB ODM};
  \node[comp] (m4) at (2.8,-1.7) {\textbf{ModuloNotificaciones}\\[1pt]Socket.io Gateway};
  \node[comp] (m5) at (5.6,-1.7) {\textbf{ModuloSeguridad}\\[1pt]Auditoría / Bloqueo};

  % Conexiones superiores
  \draw[arrow] (0,5.15) -- node[right, font=\tiny] {Interacción UI} (front.north);
  \draw[arrow] (front.south) -- node[right, font=\tiny] {Petición HTTP} (api.north);
  \draw[arrow] (api.south) -- node[right, font=\tiny] {Evento Disparado} (bus.north);

  % Modulo 1: Monitoreo
  \draw[arrow] (m1.north) to[bend left=16] node[left, font=\tiny, sloped] {Evento MetricaRecolectada} (bus.west);
  \draw[arrow] (m1.east) to[out=25,in=-25,looseness=3] node[right, font=\tiny] {Sondeo ICMP/ARP} (m1.east);

  % Modulo 2: Alertas
  \draw[arrow] (bus.south west) to[bend right=15] node[left, font=\tiny] {Consume Metrica} (m2.north);
  \draw[arrow] (m2.north) to[bend right=15] node[right, font=\tiny] {Evento AlertaGenerada} (bus.south west);
  \draw[arrow] (m2.east) to[out=25,in=-25,looseness=3] node[right, font=\tiny] {Evalúa umbral} (m2.east);

  % Modulo 3: Persistencia
  \draw[arrow] (bus.south) -- node[left, font=\tiny] {Consume Metrica/Alerta} (m3.north);
  \draw[arrow] (m3.east) to[out=25,in=-25,looseness=3] node[right, font=\tiny] {Persiste BSON} (m3.east);

  % Modulo 4: Notificaciones
  \draw[arrow] (bus.south east) to[bend left=15] node[right, font=\tiny] {Consume Alerta/Estado} (m4.north);
  \draw[arrow] (m4.east) to[out=25,in=-25,looseness=3] node[right, font=\tiny] {Emite WebSocket} (m4.east);
  \draw[arrow, draw=espochBlue, thick, dashed] (m4.north) to[bend right=38] node[right, font=\tiny] {WebSocket push} (front.east);

  % Modulo 5: Seguridad
  \draw[arrow] (bus.east) to[bend left=16] node[right, font=\tiny, sloped] {Consume IntentoFallido} (m5.north);
  \draw[arrow] (m5.east) to[out=25,in=-25,looseness=3] node[right, font=\tiny] {Bloquea cuenta} (m5.east);

\end{tikzpicture}
\caption{Diagrama de componentes y flujo de eventos del sistema}\label{fig:diagrama_eventos}
\end{figure}

\section{Catálogo Estandarizado de Eventos}
El bus de eventos opera mediante constantes inmutables para impedir el uso de cadenas mágicas y asegurar tipado estricto en el backend:

\begin{table}[htbp]
\centering
\caption{Catálogo de eventos del bus interno}\label{tab:catalogo_eventos}
\small
\begin{tabularx}{\textwidth}{>{\raggedright\arraybackslash\ttfamily\footnotesize}p{4.1cm} >{\raggedright\arraybackslash\small}p{2.5cm} >{\raggedright\arraybackslash\small}p{3.3cm} >{\raggedright\arraybackslash\small}X}
\toprule
\textbf{\textnormal{\textsf{\textbf{Evento Tipado}}}} & \textbf{\textnormal{\textsf{\textbf{Productor}}}} & \textbf{\textnormal{\textsf{\textbf{Consumidor(es)}}}} & \textbf{\textnormal{\textsf{\textbf{Acción Desencadenada}}}} \\
\midrule
NETWORK\_PING\_RESULT & ModuloMonitoreo & ModuloAlertas, ModuloPersistencia & Evalúa umbrales de latencia y persiste la telemetría en MongoDB. \\
\midrule
DEVICE\_STATUS\_CHANGED & ModuloMonitoreo & ModuloAlertas, ModuloNotificaciones & Actualiza estado (\texttt{online}, \texttt{warning}, \texttt{offline}) y emite por WebSocket. \\
\midrule
DEVICE\_DISCOVERED & ModuloMonitoreo & ModuloNotificaciones & Notifica descubrimiento de un nuevo host no inventariado en el segmento LAN. \\
\midrule
ALERT\_TRIGGERED & ModuloAlertas & ModuloNotificaciones, ModuloPersistencia & Emite alerta visual y acústica al operador y crea el ticket en base de datos. \\
\midrule
AUTH\_LOGIN\_FAILED & ModuloControlador & ModuloSeguridad & Incrementa contador de fallos; tras 5 intentos ejecuta bloqueo atómico temporal. \\
\midrule
TOPOLOGY\_LINK\_CHANGED & ModuloControlador & ModuloNotificaciones & Actualiza aristas del grafo de red en tiempo real en la vista de topología. \\
\bottomrule
\end{tabularx}
\end{table}

\newpage
\chapter{Desglose Técnico del Stack Tecnológico}

\section{Componentes Seleccionados y Funciones}
A continuación se detalla la función técnica específica de cada tecnología seleccionada:

\begin{itemize}
  \item \textbf{React 18.2+ con Vite 5:} Interfaz reactiva concurrente. Actualiza el mapa de red y paneles métricos en tiempo real sin recargar la página.
  \item \textbf{Socket.io (v4.8 Backend / v4.7 Frontend):} Canal WebSocket persistente para transmisión bidireccional en tiempo real con reconexión automática y latencia inferior a 50 ms.
  \item \textbf{React Flow (@xyflow/react v12):} Motor de visualización para grafos interactivos de topología de red con nodos personalizados, aristas dinámicas, zoom y arrastre.
  \item \textbf{Node.js v20 LTS y Express v4.21+:} Entorno asíncrono no bloqueante con microframework para gestión de rutas REST, validación Zod y middlewares de seguridad (Helmet, CORS).
  \item \textbf{Bus de Eventos Interno (\texttt{EventEmitter}):} Despachador de mensajes en memoria que desacopla la captura de métricas, el motor de alertas y la persistencia.
  \item \textbf{Mecanismos de Red Nativos (\texttt{arp-scan} y \texttt{ping}):} Ejecución eficiente mediante procesos del kernel de Linux con pool de concurrencia (\texttt{p-limit} a 15 sockets simultáneos).
  \item \textbf{MongoDB 7.0 con Mongoose 8.9+:} Base de datos NoSQL documental con índices TTL para purga automática de sesiones y métricas históricas a 30 y 90 días.
  \item \textbf{Nginx 1.22+ y Systemd en Debian 12:} Servidor web inverso con terminación SSL y proxy WebSocket, orquestado como servicio nativo con política \texttt{Restart=always}.
\end{itemize}

\begin{figure}[htbp]
\centering
\begin{tikzpicture}[
  scale=0.88, every node/.style={transform shape},
  layer/.style={rectangle, draw=black, thick, fill=techGray, rounded corners=3pt, minimum height=1.2cm, minimum width=13.0cm, align=left, font=\small}
]

  % Capas tecnologicas
  \node[layer, fill=white, draw=espochBlue] (l1) at (0,3.6) {\textbf{Capa de Presentación Web:}\newline React 18 $\cdot$ Vite $\cdot$ React Flow (@xyflow) $\cdot$ Tailwind CSS $\cdot$ Socket.io Client};

  \node[layer, fill=white, draw=espochBlue] (l2) at (0,1.8) {\textbf{Capa de Negocio y Eventos:}\newline Node.js v20 $\cdot$ Express $\cdot$ Socket.io Server $\cdot$ EventEmitter $\cdot$ Zod $\cdot$ JWT Dual};

  \node[layer, fill=white, draw=espochRed] (l3) at (0,0.0) {\textbf{Capa de Monitoreo de Red:}\newline Daemon Linux $\cdot$ \texttt{arp-scan} (Capa 2) $\cdot$ ICMP Ping (Capa 3) $\cdot$ \texttt{p-limit} Concurrencia};

  \node[layer, fill=white, draw=espochBlue] (l4) at (0,-1.8) {\textbf{Capa de Persistencia e Infraestructura:}\newline MongoDB 7.0 (Índices TTL) $\cdot$ Mongoose $\cdot$ Debian 12 $\cdot$ Nginx $\cdot$ Systemd};

  % Flechas entre capas
  \draw[-{Stealth[scale=1.0]}, thick] (l1.south) -- (l2.north);
  \draw[-{Stealth[scale=1.0]}, thick] (l2.south) -- (l3.north);
  \draw[-{Stealth[scale=1.0]}, thick] (l3.south) -- (l4.north);

\end{tikzpicture}
\caption{Capas técnicas y distribución de responsabilidades del stack}\label{fig:stack_layers}
\end{figure}

\newpage
\chapter{Matriz de Justificación Técnica y Alternativas Descartadas}

\begin{table}[htbp]
\centering
\caption{Matriz comparativa de arquitectura, tiempo real y desarrollo web}\label{tab:comparativa_software}
\small
\begin{tabularx}{\textwidth}{p{2.6cm} p{2.8cm} p{3.8cm} X}
\toprule
\textbf{Capa / Componente} & \textbf{Tecnología Elegida} & \textbf{Alternativas Evaluadas} & \textbf{Criterio y Justificación de Decisión} \\
\midrule
\textbf{Patrón Arquitectural} & \textbf{Monolito Modular Dirigido por Eventos (EDA)} & Microservicios distribuidos, Monolito síncrono N-capas & Microservicios introduce sobrecarga de latencia de red entre servicios y complejidad innecesaria para un equipo de 4 personas. El monolito síncrono no puede responder en tiempo real a eventos de red. EDA modular ofrece el equilibrio óptimo entre desacoplamiento y simplicidad operativa. \\
\midrule
\textbf{Comunicación en Tiempo Real} & \textbf{WebSockets (Socket.io)} & HTTP Polling periódico, Server-Sent Events (SSE) & HTTP Polling satura la red con peticiones vacías cada pocos segundos y genera latencia inaceptable. SSE solo permite flujo unidireccional servidor-cliente. Socket.io ofrece transporte bidireccional, reconexión inteligente, salas de aislamiento y latencia inferior a 50 ms. \\
\midrule
\textbf{Frontend Web} & \textbf{React 18 + Vite} & Angular, Vue.js, Plantillas Server-Side (EJS/Django) & React cuenta con el ecosistema de visualización de grafos de red más maduro (@xyflow/react). Vite ofrece compilaciones ultrarrápidas y menor consumo de memoria en desarrollo que Webpack. \\
\midrule
\textbf{Backend \& API} & \textbf{Node.js + Express} & Python (Django/FastAPI), Java (Spring Boot) & El modelo de I/O asíncrono no bloqueante de Node.js es superior al manejar cientos de comprobaciones concurrentes de red sin levantar hilos pesados del sistema operativo. \\
\bottomrule
\end{tabularx}
\end{table}

\begin{table}[htbp]
\centering
\caption{Matriz comparativa de persistencia, monitoreo de red e infraestructura}\label{tab:comparativa_infra}
\small
\begin{tabularx}{\textwidth}{p{2.6cm} p{2.8cm} p{3.8cm} X}
\toprule
\textbf{Capa / Componente} & \textbf{Tecnología Elegida} & \textbf{Alternativas Evaluadas} & \textbf{Criterio y Justificación de Decisión} \\
\midrule
\textbf{Base de Datos} & \textbf{MongoDB 7.0} & PostgreSQL, MySQL, SQLite & La estructura de los dispositivos de red es polimórfica (un router tiene tablas de rutas; un switch puertos y VLANs; una PC no). Los documentos BSON modelan esta jerarquía sin tablas relacionales artificiales. Dispone de índices TTL nativos. \\
\midrule
\textbf{Daemon de Red} & \textbf{Node.js Daemon + arp-scan/ping} & Python Scapy, C/C++ raw sockets & `arp-scan` y `ping` son binarios del sistema altamente optimizados por el kernel de Linux. Integrarlos con Node.js proporciona equilibrio entre velocidad de ejecución en bajo nivel y facilidad de lógica de negocio. \\
\midrule
\textbf{Despliegue} & \textbf{Systemd + Nginx en Debian} & Docker Compose en VM, Kubernetes & En la máquina física del laboratorio asignada para el proyecto, capas pesadas de virtualización consumen memoria RAM innecesaria. La ejecución como servicio nativo Systemd garantiza el 100\% del rendimiento del hardware. \\
\bottomrule
\end{tabularx}
\end{table}

\end{document}
